%%%%%%%%%%  scan idx 370  |  folio 357  %%%%%%%%%%
\SGAPrintPage{357}%
Since \(A\) is of finite grading, the graded rings associated with the \(\lambda\)-ring filtration and the degree filtration on \(\Ch A\) are identical (6.6.7). We then define \(u:A\longrightarrow \Gr_{\deg}(\Ch A)\) by the identity on \(A_0\) and by:
\[
  \forall\, i\ge1,\quad\forall\, a_i\in A_i,\qquad u(a_i)=\cl_{\Gr^i}\left(0,1+(-1)^{i-1}(i-1)!a_i\right),
\]
and extend by additivity. By relation (6.5.3), this is a homomorphism of graded rings, and it is easy to see that it is compatible with the operations of \(K\). The inverse morphism \(u^{-1}\) can be defined as follows: \(\forall\, i\ge1\), \(\forall\,\bar\varphi\in\Gr^i(\Ch A)\), the \(i\)-th coefficient \(a_i\) of \(\varphi\) does not depend on the way in which \(\bar\varphi\) has been lifted to \((\Ch A)_i\), and we put:
\[
  u^{-1}(\bar\varphi)=\frac{(-1)^{i-1}}{(i-1)!}a_i.
\]
The homomorphism \(u\) is therefore an isomorphism.

The isomorphism \(\mathord{\wedge}\to\Ch(\Gr\mathord{\wedge})\) is defined by the completed Chern class \(\widetilde c\). To verify that it does define an isomorphism, it suffices to verify that the homomorphism obtained by passing to the associated graded rings is one, since the filtrations are discrete. We therefore obtain:
\[
  \Gr\mathord{\wedge}
   \xrightarrow{\Gr(\widetilde c)}
  \Gr_{\lambda}(\Ch(\Gr\mathord{\wedge}))
   \simeq
  \Gr_{\deg}(\Ch(\Gr\mathord{\wedge}))
   \xrightarrow{u^{-1}}
  \Gr\mathord{\wedge} .
\]
It suffices to show that the composite is the identity. Since the homomorphisms are homomorphisms of graded \(K\)-algebras, it suffices, by the definition of the \(\lambda\)-ring filtration on \(\mathord{\wedge}\), to show that it is the identity on the \(\overline{\gamma^i(x)}\), with \(\epsilon(x)=0\). We then obtain:
\[
  \Gr(\widetilde c)\bigl(\overline{\gamma^i(x)}\bigr)
  =
  \cl_{\Gr^i}\bigl(\gamma^i(\widetilde c(x))\bigr).
\]
Now \(\gamma^i(\widetilde c(x))=(0,1+(-1)^{i-1}(i-1)!c^i(x)+\ldots)\), by (6.6.4); hence:
\[
  \cl_{\Gr^i}\bigl(\gamma^i(\widetilde c(x))\bigr)=u(c^i(x)),
\]
%%%%%%%%%%  scan idx 371  |  folio 358  %%%%%%%%%%
\SGAPrintPage{358}%
which shows that \(u^{-1}\circ\Gr(\widetilde c)=\Id\).

\sgaStableTarget{sga6.e05.cor.6-12}\textbf{Corollary 6.12.} \uline{If} \(K\) \uline{is of characteristic zero}, \uline{and} \(\mathord{\wedge}\) \uline{is a} \(K\)-\(\lambda\)-\uline{algebra} \uline{with discrete filtration}, \uline{then the completed Chern class}:
\[
  \widetilde c:\mathord{\wedge}\longrightarrow \Ch(\Gr\mathord{\wedge})
\]
\uline{is an isomorphism}.

The proof is contained in that of 6.11.

\sgaStableTarget{sga6.e05.def.6-13}\textbf{Definition 6.13.} \uline{Let} \(A\) \uline{be a graded} \(K\)-\uline{algebra}, \uline{and} \(\mathord{\wedge}\) \uline{an} \uline{augmented} \(K\)-\(\lambda\)-\uline{algebra}; \uline{we call a Chern theory for} \(\mathord{\wedge}\) \uline{with values in} \(A\) \uline{any} \(\lambda\)-\uline{homomorphism}: \(\mathord{\wedge}\longrightarrow \Ch(A)\), \uline{compatible with the augmentations}.

We define a functor on the category of graded \(K\)-algebras, with values in the category of sets, by associating with every graded \(K\)-algebra \(A\) the set of Chern theories for \(\mathord{\wedge}\) with values in \(A\), namely \(\Hom_\lambda(\mathord{\wedge},\Ch(A))\).

\sgaStableTarget{sga6.e05.prop.6-14}\textbf{Proposition 6.14.} \uline{The functor defined in} 6.13 \uline{is representable}. \uline{If} \(K\) \uline{is} \uline{of characteristic zero}, \uline{and} \(\mathord{\wedge}\) \uline{has discrete filtration}, \uline{it is represented} \uline{by the graded ring associated with} \(\mathord{\wedge}\).

Let us first prove the first assertion. We consider the graded \(K\)-algebra generated by generators \(c^i(x)\), with \(i\ge1\), \(x\in\Fil^1\mathord{\wedge}\), where \(c^i(x)\) has degree \(i\), and take its quotient by the following relations:
\[
\begin{array}{ll}
\text{i)}&
 c^k(x+y)=\displaystyle\sum_{i+j=k}c^i(x)c^j(y),
 \quad\text{for } k\ge1,\ x,y\in\Fil^1\mathord{\wedge}\ ;\\[0.8em]
\text{ii)}&
 c^k(xy)=P''_k(c^i(x),c^j(y))
 \quad(\text{cf. }6.1),
 \quad\text{for } k\ge1,\ x,y\in\Fil^1\mathord{\wedge}\ ;\\[0.4em]
\text{iii)}&
 c^k(ax)=R_{k,a}(c^i(x)) \quad(\text{cf. }6.1),
 \quad\text{for } k\ge1,\ a\in K,\ x\in\Fil^1\mathord{\wedge}\ ;\\[0.4em]
\text{iv)}&
 c^k(\gamma^j(x))=Q''_{j,k}(c^i(x))
 \quad(\text{cf. }6.1),
 \quad\text{for } j,k\ge1,\ x\in\Fil^1\mathord{\wedge},
\end{array}
\]

%%%%%%%%%%  scan idx 372  |  folio 359  %%%%%%%%%%
\SGAPrintPage{359}%
These relations being isobaric, the quotient is a graded \(K\)-algebra \(\Omega\). We define a morphism \(\mathord{\wedge}\xrightarrow{\widetilde c_0}\Ch(\Omega)\) by:
\[
  x\rightsquigarrow
  \bigl(\epsilon(x),1+c^1(x-\epsilon(x))+\ldots\bigr).
\]
The imposed relations show that this is a Chern theory for \(\mathord{\wedge}\) with values in \(\Omega\). If in addition we have a Chern theory \(f\), with values in a graded \(K\)-algebra \(A\), there is a unique homomorphism: \(\Omega\to A\), such that \(f\) factors through \(\Ch(\Omega)\); it is defined by sending, for every \(i\ge1\) and every \(x\in\Fil^1\mathord{\wedge}\), the element \(c^i(x)\) to the \(i\)-th coefficient of \(f(x)\) (which is well defined since \(f\) is a \(\lambda\)-homomorphism).

Now suppose that \(K\) is of characteristic zero and \(\mathord{\wedge}\) has discrete filtration; let \(n\) be such that \(\Fil^n\mathord{\wedge}=0\). Then the graded algebra \(\Omega\) is zero in degrees \(\ge n\). To prove this, we must show that for \(x_1,\ldots,x_p\in\Fil^1\mathord{\wedge}\) and \(k_1+\ldots+k_p\ge n\), we have:
\[
  c^{k_1}(x_1)\mathinner{....}c^{k_p}(x_p)=0.
\]
Now, by hypothesis, we have:
\[
  \gamma^{k_1}(x_1)\mathinner{....}\gamma^{k_p}(x_p)=0\ ;
\]
Applying the \(\lambda\)-homomorphism \(\widetilde c_0\), we obtain:
\[
  \gamma^{k_1}(\widetilde c_0(x_1))\mathinner{....}
  \gamma^{k_p}(\widetilde c_0(x_p))=(0,1).
\]
Now, by (6.6.4) and (6.5.3), this implies:
\[
  (k_1+\ldots+k_p)!\,
  c^{k_1}(x_1)\ldots c^{k_p}(x_p)=0.
\]
Since \(K\) is of characteristic zero, we have the desired result.

The universal property of \(\Omega\), applied to the canonical homomorphism \(\widetilde c\), gives a commutative diagram:
%%%%%%%%%%  scan idx 373  |  folio 360  %%%%%%%%%%
\SGAPrintPage{360}%
\[
\begin{tikzcd}[column sep=large,row sep=large]
\mathord{\wedge} \arrow[r,"\widetilde c"] \arrow[d,"\widetilde c_0"'] &
\Ch(\Gr\mathord{\wedge})\\
\Ch(\Omega) \arrow[ur] &
\end{tikzcd}
\]
By Proposition 6.11, the homomorphism: \(\Omega\to\Gr\mathord{\wedge}\) is an isomorphism if and only if the homomorphism: \(\Ch(\Omega)\to\Ch(\Gr\mathord{\wedge})\) is one. Since it is visibly surjective, the latter is also surjective, and it therefore suffices to show that \(\Ch(\Omega)\to\Ch(\Gr\mathord{\wedge})\) is injective. By functoriality of \(\Ch\), we obtain the diagram:
\[
\begin{tikzcd}[column sep=large,row sep=large]
\mathord{\wedge} \arrow[r,"\widetilde c_0"] \arrow[d,"\widetilde c"'] &
\Ch(\Omega) \arrow[d,"\widetilde c"]\\
\Ch(\Gr\mathord{\wedge}) \arrow[r,"\Ch(\widetilde c_0)"] &
\Ch\bigl(\Gr_\lambda(\Ch(\Omega))\bigr);
\end{tikzcd}
\]
since \(\Omega\) has bounded degrees, the right vertical arrow is an isomorphism (6.12), and using 6.11, we therefore obtain a morphism: \(\Gr\mathord{\wedge}\longrightarrow\Omega\), the composite of the two being the identity by the universal property of \(\Omega\).

\sgaStableTarget{sga6.e05.sec.7}\section*{APPENDIX: The Adams operations \(\psi^k\)}

We follow here a presentation due to J.P. Serre.

\sgaStableTarget{sga6.e05.para.7-1}7.1. Let \(K\) be a pre-\(\lambda\)-ring. For every \(x\in K\), we denote by \(\frac{d}{dt}(\lambda_t(x))\) the derivative with respect to \(t\) of \(\lambda_t(x)\). We then define elements of \(K\), denoted by \(\psi^k(x)\), by putting:
\sgaStableTarget{sga6.e05.eq.7-1-1}\begin{equation}\tag{7.1.1}
  \frac{d}{dt}(\lambda_t(x))/\lambda_t(x)
  =
  \sum_{k=1}^{\infty}(-1)^{k-1}\psi^k(x)t^{k-1}.
\end{equation}

The maps \(\psi^k:K\to K\) are additive; indeed, we have:
\[
\begin{aligned}
 \frac{d}{dt}(\lambda_t(x+y))/\lambda_t(x+y)
 &=
 \frac{d}{dt}(\lambda_t(x)\lambda_t(y))/(\lambda_t(x)\lambda_t(y))\\
 &=
 \frac{d}{dt}(\lambda_t(x))/\lambda_t(x)
 +\frac{d}{dt}(\lambda_t(y))/\lambda_t(y)\ ;
\end{aligned}
\]
%%%%%%%%%%  scan idx 374  |  folio 361  %%%%%%%%%%
\SGAPrintPage{361}%
and therefore, for every \(k\):
\[
  \psi^k(x+y)=\psi^k(x)+\psi^k(y).
\]

By (7.1.1), the \(\psi^k\) are isobaric polynomials of weight \(k\) in the \(\lambda^i\), with \(i=1,\ldots,k\). For example, let us give the values of \(\psi^1\) and \(\psi^2\):
\[
  \psi^1(x)=\lambda^1(x)=x,
  \qquad
  \psi^2(x)=x^2-2\lambda^2(x).
\]

\sgaStableTarget{sga6.e05.para.7-2}7.2. Let \(K\) be a ring of characteristic \(0\); suppose given additive maps \(\psi^k\) from \(K\) to \(K\), such that \(\psi^1=\operatorname{Id}_K\). We then define a pre-\(\lambda\)-structure on \(K\) by putting:
\sgaStableTarget{sga6.e05.eq.7-2-1}\begin{equation}\tag{7.2.1}
  \lambda_t(x)=
  \exp\left(\sum_{k=1}^{\infty}(-1)^{k-1}\psi^k(x)t^k/k\right).
\end{equation}
Multiplicativity follows from the additivity of the \(\psi^k\). There is therefore an equivalence between giving a pre-\(\lambda\)-structure on \(K\) and giving the corresponding operations \(\psi^k\).

\sgaStableTarget{sga6.e05.para.7-3}7.3. It is clear from formula (7.1.1) that the \(\psi^k\) commute with \(\lambda\)-homomorphisms. If \(K,K'\) are pre-\(\lambda\)-rings of characteristic zero, it is necessary and sufficient that a homomorphism \(f:K\to K'\) commute with the \(\psi^k\) for it to be a \(\lambda\)-homomorphism: this follows from 7.2, the \(\lambda^i\) being expressed as polynomials in the \(\psi^k\).

\sgaStableTarget{sga6.e05.ex.7-4}\textbf{Examples 7.4.} i) For \(K=\mathbb Z\), endowed with its \(\lambda\)-ring structure, we have:
\[
  \frac{d}{dt}(\lambda_t(n))/\lambda_t(n)
  =
  n(1+t)^{n-1}/(1+t)^n
  =
  n.\sum_{k=1}(-1)^{k-1}t^{k-1}.
\]
Consequently, the \(\psi^k\) are all equal to the identity map of \(\mathbb Z\).

ii) Let us work in the category \(\mathcal C_0\) of \(\mathbb Z\)-algebras (1.8). The
